\documentclass[11pt,a4paper]{article}

% ── Packages ──────────────────────────────────────────────────────────────
\usepackage[utf8]{inputenc}
\usepackage[T1]{fontenc}
\usepackage{lmodern}
\usepackage[margin=2.5cm]{geometry}
\usepackage{amsmath,amssymb}
\usepackage{graphicx}
\usepackage{booktabs}
\usepackage{hyperref}
\usepackage{cleveref}
\usepackage{natbib}
\usepackage{xcolor}
\usepackage{enumitem}
\usepackage{microtype}

\hypersetup{
  colorlinks=true,
  linkcolor=blue!60!black,
  citecolor=blue!60!black,
  urlcolor=blue!60!black
}

% ── Metadata ──────────────────────────────────────────────────────────────
\title{\texttt{maxentcpp}: A C++ reimplementation of Maximum Entropy
       species distribution modeling for R}

\author{
  \'Angel Luis Robles Fern\'andez\thanks{Corresponding author.
    ORCID: \href{https://orcid.org/0000-0002-4741-8012}{0000-0002-4741-8012}.
    E-mail: \texttt{a.l.robles.fernandez@gmail.com}}\\[4pt]
  \small Department of Ecology and Evolutionary Biology,\\
  \small University of Kansas, Lawrence, KS, United States
}

\date{May 2026}

\begin{document}
\maketitle

% ══════════════════════════════════════════════════════════════════════════
\section{Summary}
% ══════════════════════════════════════════════════════════════════════════

\texttt{maxentcpp} is an R package that provides a high-performance
C++17 reimplementation of the Maximum Entropy (Maxent) algorithm for
species distribution modeling~(SDM).
Maxent estimates the geographic distribution of a species from
presence-only occurrence records and environmental covariates by finding
the probability distribution of maximum entropy subject to constraints
derived from the training data
\citep{Phillips2006,Phillips2004}.
Since its introduction, Maxent has become one of the most widely cited
methods in ecology and conservation biology, with over 13\,000 citations
\citep{Elith2011}.

The fitted model is a \emph{Gibbs distribution} over geographic space:
\begin{equation}\label{eq:gibbs}
  P(\mathbf{x}) \;=\;
    \frac{1}{Z}\,
    \exp\!\Bigl(\sum_{j=1}^{J} \lambda_j\, f_j(\mathbf{x})\Bigr),
\end{equation}
where $f_j$ are feature transformations of the environmental variables,
$\lambda_j$ are learned weights, and
$Z = \sum_{\mathbf{x}}\exp\!\bigl(\sum_j \lambda_j f_j(\mathbf{x})\bigr)$
is the normalizing constant (partition function).
The optimizer finds the $\lambda_j$ that maximize the regularized
log-likelihood via sequential coordinate ascent with
$\ell_1$ (lasso) penalties \citep{Phillips2006,Elith2011}.

The original Maxent software is a Java desktop application
\citep{Phillips2017}.
\texttt{maxentcpp} translates the complete Maxent~3.4.4 algorithm into
C++17 with R~bindings via Rcpp \citep{Eddelbuettel2011} and
RcppEigen \citep{Bates2013}, eliminating the Java Runtime Environment
dependency entirely.
The package supports all five feature types (linear, quadratic, product,
hinge, and threshold), output transformations in raw, logistic, and
complementary log-log (cloglog) scales, and a streaming raster
evaluation engine that processes environmental layers block-by-block
through the \texttt{terra} package \citep{Hijmans2024}, enabling
continental-scale projections on commodity hardware without loading
entire raster stacks into memory.
Model diagnostics include AUC evaluation, permutation importance,
response curves, and Multivariate Environmental Similarity Surfaces
(MESS) for novelty detection.

% ══════════════════════════════════════════════════════════════════════════
\section{Statement of need}
% ══════════════════════════════════════════════════════════════════════════

Maxent is the \emph{de facto} standard for presence-only species
distribution modeling \citep{Merow2013}.
However, the original Java implementation presents several practical
barriers for modern ecological research workflows:

\begin{enumerate}[leftmargin=*,nosep]
  \item \textbf{Java dependency.}
    Users must install and configure a compatible JDK\@.
    Version conflicts between Java~8, 11, 17, and~21 produce silent
    failures, and the \texttt{rJava} bridge package is one of the most
    common sources of installation errors in the R~ecosystem,
    particularly on macOS and HPC clusters.
    While \texttt{conda} environments and Docker containers can manage
    Java versioning, they add deployment complexity and are not standard
    practice in many ecology research groups.

  \item \textbf{File-based I/O.}
    Data must be serialized to disk as SWD or ASCII grid files before
    the Java process can read them, and results must be parsed back
    from output files.
    There is no in-memory bridge between R data structures and the Java
    optimizer.

  \item \textbf{Scalability.}
    The Java implementation is single-threaded and GUI-centric.
    In multi-species, multi-scenario workflows (e.g., 500
    species $\times$ 4 climate scenarios), file~I/O overhead
    accumulates: each species run requires writing SWD files, invoking
    the JVM, and parsing output files.
    A~native in-memory API eliminates this per-call overhead entirely.

  \item \textbf{Maintenance.}
    The Java Maxent codebase has remained essentially unchanged since
    version~3.4.4 (released November 2020), with only two active
    contributors \citep{Phillips2017}.
\end{enumerate}

\texttt{maxentcpp} addresses these barriers by providing a native C++/R
implementation that installs with a single
\texttt{install.packages("maxentcpp")} call, operates entirely in memory
through R objects, and integrates seamlessly with the \texttt{terra}
spatial data ecosystem.
The target audience is ecologists, conservation biologists, and
biogeographers who use Maxent in reproducible, script-based
workflows---from individual species analyses to large-scale biodiversity
assessments.

% ══════════════════════════════════════════════════════════════════════════
\section{Legacy software modernization in ecology}
% ══════════════════════════════════════════════════════════════════════════

Computational ecology relies on a small number of foundational software
tools, many of which were written over a decade ago and have not been
modernized for current computing environments.
Reimplementing legacy scientific software in modern languages is uncommon
in ecology but has important precedents in adjacent environmental
sciences.

Circuitscape, the standard tool for landscape connectivity modeling, was
rewritten from Python to Julia, achieving order-of-magnitude speedups
while preserving result equivalence \citep{Hall2021}.
The LANDIS-II forest landscape model was re-engineered from monolithic~C
into a modular C\# framework with automated testing and version control
\citep{Scheller2009}.
Most recently, the WaterGAP global hydrological model was reprogrammed
from Fortran to Python with systematic cross-validation against the
original \citep{Nyenah2025}.
Closer to species distribution modeling, \texttt{maxnet}
\citep{Phillips2024maxnet} represents the original MaxEnt author's own
reimplementation using~R and \texttt{glmnet}---preserving the same
statistical model but employing a different optimization algorithm
(coordinate descent on the full feature matrix rather than sequential
coordinate ascent).

These precedents share common motivations: eliminating obsolete runtime
dependencies, enabling integration into modern scripted workflows,
improving maintainability, and opening the algorithm to inspection and
extension.
\texttt{maxentcpp} follows this lineage but goes further by
\textbf{preserving the original optimization algorithm} (sequential
coordinate ascent with $\ell_1$-penalized Newton steps) rather than
substituting an approximate equivalent, and by \textbf{validating
per-iteration numerical equivalence} rather than only comparing
final-state outputs.

To our knowledge, \texttt{maxentcpp} is the first project in ecology to
systematically reimplement a widely used SDM algorithm in a compiled
language with per-iteration numerical validation against the original
implementation, combining faithful algorithm porting with quantitative
fidelity testing as a companion package (\texttt{maxentcppCompTest}).

% ══════════════════════════════════════════════════════════════════════════
\section{State of the field}
% ══════════════════════════════════════════════════════════════════════════

Several R packages provide access to Maxent-family models.
\textbf{dismo} \citep{Hijmans2023} is the most widely used R interface
to Java Maxent; it wraps \texttt{maxent.jar} via \texttt{rJava},
inheriting the JDK dependency and file~I/O overhead, and is currently in
maintenance mode following the transition from \texttt{raster} to
\texttt{terra}.
\textbf{maxnet} \citep{Phillips2024maxnet} is a pure-R implementation
that uses \texttt{glmnet} for elastic-net regularization; while it
avoids the Java dependency, it employs a different optimization
algorithm (elastic net via coordinate descent on the full feature
matrix) rather than the sequential coordinate-ascent optimizer of the
original Maxent, which can produce numerically different results,
particularly for threshold and hinge features.
\textbf{ENMeval} \citep{Kass2021} is a model tuning and evaluation
framework that wraps either \texttt{dismo::maxent()} or
\texttt{maxnet::maxnet()}; it does not implement the Maxent algorithm
itself.
\textbf{kuenm} \citep{Cobos2019} is a calibration and evaluation
toolkit that relies on \texttt{dismo} or direct Java Maxent calls.

\subsection{Empirical comparison: maxentcpp vs maxnet}

To quantify the practical differences between \texttt{maxentcpp} and
\texttt{maxnet}, we compared both packages on a virtual species
\citep{Leroy2016} with known true habitat suitability, constructed from
the environmental layers bundled with \texttt{maxentcpp} (Annual Mean
Temperature and Annual Precipitation over Central America, 2,371
non-NA cells).
The virtual species has a Gaussian niche centered at 20\,\textdegree{}C
and 1,500\,mm, and 100 presence records were sampled proportionally to
the true suitability surface (see the companion vignette
\emph{Virtual Species Comparison}\footnote{%
  \url{https://alrobles.github.io/maxentcpp/articles/virtual_species_comparison.html}}
for full reproducible code).

Both packages were fitted with linear, quadratic, and hinge features.
The results demonstrate high ecological equivalence:

\begin{table}[htbp]
\centering
\caption{Ecological equivalence metrics: \texttt{maxentcpp} vs
         \texttt{maxnet} on a virtual species.}
\label{tab:equivalence}
\begin{tabular}{@{}lr@{}}
\toprule
Metric & Value \\
\midrule
Schoener's $D$ & 0.93 \\
Spearman $\rho$ (rank correlation) & 0.95 \\
Pearson $r$ & 0.97 \\
Mean $|\Delta\text{cloglog}|$ & 0.053 \\
Max $|\Delta\text{cloglog}|$ & 0.43 \\
\bottomrule
\end{tabular}
\end{table}

These results confirm that \texttt{maxentcpp} and \texttt{maxnet}
produce \textbf{ecologically equivalent predictions} for typical
modeling scenarios (Schoener's $D > 0.9$ is conventionally interpreted
as high niche overlap).
The largest differences occur in the tails of the suitability
distribution, where the two optimization algorithms
(\texttt{maxentcpp}: sequential coordinate ascent; \texttt{maxnet}:
elastic-net coordinate descent) regularize differently.
Both packages recover the true niche with comparable accuracy
(Spearman $\rho > 0.87$ vs true surface).

The key practical scenarios where a user must choose \texttt{maxentcpp}
over \texttt{maxnet} are:
(1)~when exact reproduction of Java Maxent results is required
(e.g., replicating published analyses),
(2)~when lambda file interoperability with Java Maxent is needed, and
(3)~when streaming raster projection onto grids larger than available
RAM is required.

\subsection{Performance benchmarks}

Training time for the bundled \emph{Abeillia abeillei} dataset
(73~presences, 2,371~background, linear + quadratic + hinge features):

\begin{table}[htbp]
\centering
\caption{Training time comparison.}
\label{tab:benchmark}
\begin{tabular}{@{}lr@{}}
\toprule
Package & Median time per fit \\
\midrule
\texttt{maxentcpp} & $\sim$820\,ms \\
\texttt{maxnet}    & $\sim$530\,ms \\
\bottomrule
\end{tabular}
\end{table}

\texttt{maxnet} is faster for small datasets because \texttt{glmnet}'s
coordinate descent is highly optimized for dense feature matrices.
However, \texttt{maxentcpp}'s streaming evaluation avoids materializing
the full prediction matrix, giving it an advantage for projection onto
large rasters where \texttt{dismo}/Java Maxent would require disk-based
intermediates.

\texttt{maxentcpp} was built as a new package rather than contributing
to existing projects for three reasons.
First, a faithful port of the original Java optimizer (sequential
coordinate ascent with $\ell_1$-penalized Newton steps) requires a C++
engine that cannot be expressed through \texttt{glmnet}'s API;
contributing this to \texttt{maxnet} would change its core algorithm.
Second, \texttt{dismo}'s architecture is tightly coupled to
\texttt{rJava} and \texttt{raster}; the native C++ approach eliminates
this dependency chain entirely.
Third, the header-based C++ library design of \texttt{maxentcpp} enables
reuse in other packages or standalone applications beyond~R---something
not achievable through modifications to existing R-only packages.

% ══════════════════════════════════════════════════════════════════════════
\section{Software design}
% ══════════════════════════════════════════════════════════════════════════

\subsection{Architecture}

\texttt{maxentcpp} is organized in three layers (C++ core, Rcpp bridge,
R~interface), with the C++ core further split into algorithmic and
infrastructure sublayers:

\begin{table}[htbp]
\centering
\caption{Architecture of \texttt{maxentcpp}: three layers with the C++
         core divided into algorithmic and infrastructure sublayers.}
\label{tab:arch}
\begin{tabular}{@{}llr@{}}
\toprule
Layer & Key components & Lines \\
\midrule
C++ core (algorithmic)
  & \texttt{Sequential}, \texttt{FeaturedSpace}, five feature types
  & $\sim$4\,500 \\
C++ core (I/O, diagnostics)
  & \texttt{BackgroundProvider}, grid I/O, CSV, MESS, response curves
  & $\sim$2\,400 \\
Rcpp bridge
  & \texttt{rcpp\_*.cpp} external-pointer bindings \citep{Eddelbuettel2013}
  & $\sim$3\,300 \\
R interface
  & \texttt{maxent\_run()}, feature generators, projection, evaluation,
    diagnostics
  & $\sim$2\,500 \\
\bottomrule
\end{tabular}
\end{table}

The C++ core uses Eigen \citep{Guennebaud2010} for dense linear algebra.
Of the $\sim$6\,900 C++ lines, approximately 4\,500 implement the core
algorithmic logic (optimizer, features, density) while the remainder
handles I/O, Rcpp bindings, and diagnostics.
The high-level \texttt{maxent\_run()} function provides a one-call
workflow mirroring the Java GUI experience, while lower-level functions
(\texttt{maxent\_generate\_features()}, \texttt{maxent\_featured\_space()},
\texttt{maxent\_fit()}, \texttt{maxent\_project\_cloglog()}) offer full
control over each modeling step.

\subsection{Optimization algorithm}

The \texttt{Sequential} optimizer is a direct port of
\texttt{density.Sequential} in Java Maxent~3.4.4.
It minimizes the $\ell_1$-regularized negative log-likelihood:
\begin{equation}\label{eq:objective}
  \mathcal{L}(\boldsymbol{\lambda})
  = -\frac{1}{m}\sum_{i=1}^{m} \sum_{j=1}^{J}
      \lambda_j\, f_j(\mathbf{x}_i)
    + \log Z(\boldsymbol{\lambda})
    + \sum_{j=1}^{J} \beta_j\,|\lambda_j|,
\end{equation}
where $m$ is the number of presence records,
$Z(\boldsymbol{\lambda})
= \sum_{k=1}^{N} \exp\!\bigl(\sum_j \lambda_j f_j(\mathbf{x}_k)\bigr)$
is the partition function summed over $N$ background points, and
$\beta_j = \hat\sigma_j / \sqrt{m}$ is the per-feature regularization
parameter with $\hat\sigma_j$ being the standard deviation of feature~$j$
over the presence points (floored at 0.001).

At each iteration, the optimizer selects the feature~$j^*$ with the most
negative $\Delta\mathcal{L}$ bound (the \texttt{deltaLossBound}
function), then computes a Newton step:
\begin{equation}\label{eq:newton}
  \alpha_j = -\frac{\partial \mathcal{L}/\partial \lambda_j}
    {\mathbf{u}_j^\top \mathbf{H}\, \mathbf{u}_j},
\end{equation}
where $\mathbf{u}_j$ is the $j$-th coordinate direction and
$\mathbf{u}_j^\top \mathbf{H}\, \mathbf{u}_j = \mathrm{Var}_q[f_j]$
is the variance of feature~$j$ under the current distribution~$q$.
The derivative includes the $\ell_1$ subgradient:
$\partial \mathcal{L}/\partial \lambda_j = \mathbb{E}_q[f_j] -
\mathbb{E}_{\hat{p}}[f_j] + \beta_j\,\mathrm{sign}(\lambda_j)$.
The step is damped during early iterations
($\alpha \leftarrow \alpha/50$ for iterations $<10$,
$\alpha/10$ for $<20$, $\alpha/3$ for $<50$) and falls back to
a line search (\texttt{searchAlpha}) if the Newton step does not
decrease loss.
Every 10 iterations, a parallel update applies coordinate steps to all
features simultaneously.

Convergence is tested every 20 iterations: training stops when the
relative loss improvement falls below $10^{-5}$ or 500 iterations are
reached.

\subsection{Streaming raster evaluation}

\texttt{maxentcpp} reads \texttt{terra::SpatRaster} objects
block-by-block through a \texttt{BackgroundProvider} abstraction.
Each tile is scored by the C++ engine and discarded before the next tile
is loaded, allowing projection onto raster stacks larger than available
RAM\@.
The partition function~$Z$ is accumulated across tiles by summing
unnormalized densities
$\exp(\sum_j \lambda_j f_j(\mathbf{x}_k))$ for each cell~$k$ in the
tile, then normalizing once after all tiles are processed.
This produces identical results to single-pass evaluation because the
sum is commutative and no tile-boundary effects exist---each cell is
scored independently.

\subsection{Numerical fidelity}

Every C++ class is a direct port of the corresponding Java class,
preserving the same control flow, regularization logic, and numerical
constants (\Cref{tab:porting}).
This design choice prioritizes backward compatibility: users migrating
from Java Maxent obtain identical results.
The fidelity contract is enforced by a dedicated companion package,
\texttt{maxentcppCompTest} \citep{maxentcppCompTest2025}, which runs the
C++ optimizer and the original Java \texttt{density.Sequential} on
shared test fixtures and compares per-iteration trajectories of loss,
entropy, and $\lambda$~vectors.

On controlled test fixtures (identical input data, same floating-point
representation), agreement reaches $< 10^{-14}$ relative error for
symmetric fixtures and $< 10^{-6}$ on
$\|\Delta\lambda\|_\infty$ for asymmetric fixtures at every iteration
checkpoint.
\textbf{These bounds hold under the specific conditions of the test
fixtures} (same compiler family, IEEE~754 double precision, deterministic
iteration order).
Floating-point ordering differences introduced by alternative BLAS
backends or aggressive compiler optimizations (e.g., \texttt{-ffast-math})
could widen the gap, though the sequential nature of the coordinate-ascent
algorithm limits sensitivity to parallelism-induced reordering.
Cross-platform CI (Ubuntu, macOS, Windows with GCC, Clang, and MSVC)
confirms that the test fixtures pass on all three compiler families.

\begin{table}[htbp]
\centering
\caption{Class-level porting map from Java Maxent to
         \texttt{maxentcpp}.}
\label{tab:porting}
\begin{tabular}{@{}ll@{}}
\toprule
\texttt{maxentcpp} (C++) & Java Maxent original \\
\midrule
\texttt{LinearFeature}      & \texttt{density.LinearFeature}    \\
\texttt{QuadraticFeature}   & \texttt{density.SquareFeature}    \\
\texttt{ProductFeature}     & \texttt{density.ProductFeature}   \\
\texttt{HingeFeature}       & \texttt{density.HingeFeature}     \\
\texttt{ThresholdFeature}   & \texttt{density.ThresholdFeature} \\
\texttt{FeaturedSpace}      & \texttt{density.FeaturedSpace}    \\
\texttt{Sequential::run()}  & \texttt{density.Sequential.run()} \\
\bottomrule
\end{tabular}
\end{table}

\paragraph{Lambda file compatibility.}
Trained models are serialized in the same \texttt{.lambdas} text format
used by Java Maxent~3.4.4.
Lambda files produced by either implementation can be loaded by the
other, ensuring interoperability with existing workflows and archived
models.

\subsection{Feature completeness}

The following table summarizes the current scope of \texttt{maxentcpp}
relative to Java Maxent~3.4.4 and \texttt{maxnet}:

\begin{table}[htbp]
\centering
\caption{Feature completeness: \texttt{maxentcpp} vs Java Maxent
         vs \texttt{maxnet}.}
\label{tab:features}
\small
\begin{tabular}{@{}lccc@{}}
\toprule
Feature & \texttt{maxentcpp} & Java Maxent & \texttt{maxnet} \\
\midrule
Linear features              & \checkmark & \checkmark & \checkmark \\
Quadratic features           & \checkmark & \checkmark & \checkmark \\
Product features             & \checkmark & \checkmark & \checkmark \\
Hinge features               & \checkmark & \checkmark & \checkmark \\
Threshold features           & \checkmark & \checkmark & \checkmark \\
Raw/logistic/cloglog output  & \checkmark & \checkmark & \checkmark \\
AUC evaluation               & \checkmark & \checkmark & via ENMeval \\
Permutation importance       & \checkmark & \checkmark & via ENMeval \\
Response curves              & \checkmark & \checkmark & partial \\
MESS maps                    & \checkmark & \checkmark & -- \\
Clamping / fade-by-clamping  & \checkmark & \checkmark & -- \\
Lambda file I/O              & \checkmark & \checkmark & -- \\
Streaming raster projection  & \checkmark & -- & -- \\
Bias grids                   & \checkmark & \checkmark & via \texttt{offset} \\
Categorical variables        & -- & \checkmark & \checkmark \\
Replicate runs / cross-val.  & -- & \checkmark & via ENMeval \\
Jackknife variable selection & -- & \checkmark & -- \\
Missing data handling        & -- & \checkmark & \checkmark \\
SWD-to-raster projection     & -- & \checkmark & -- \\
\bottomrule
\end{tabular}
\end{table}

Categorical variables, replicate runs, jackknife, and SWD-to-raster
projection are not yet implemented.
Users requiring these features should continue using Java Maxent or
\texttt{maxnet}.
We plan to add categorical variable support in a future release.

\subsection{Development provenance}

The translation followed a phased approach across four publicly
accessible repositories:

\begin{enumerate}[leftmargin=*,nosep]
  \item \texttt{alrobles/Maxent}---a fork of the original
    \texttt{mrmaxent/Maxent} \citep{Phillips2017} Java source code,
    where the architecture was analyzed and each Java class was mapped to
    a C++ target (193~commits).
  \item \texttt{alrobles/maxentcpp-devel}---the development repository
    where the C++ port, Rcpp bindings, R interface, tests, and
    documentation were iteratively built and refined (34~commits).
  \item \texttt{alrobles/maxentcppCompTest}---the cross-language fidelity
    test suite containing Java oracle runners, golden trajectory files,
    and automated comparison scripts (40~commits).
  \item \texttt{alrobles/maxentcpp}---the release repository with the
    clean, CRAN-ready package (47~commits).
\end{enumerate}

% ══════════════════════════════════════════════════════════════════════════
\section{Ecosystem comparison}
% ══════════════════════════════════════════════════════════════════════════

The R ecosystem contains several packages that provide access to MaxEnt
models, each employing a distinct integration strategy:

\begin{table}[htbp]
\centering
\caption{R packages providing MaxEnt access.}
\label{tab:ecosystem}
\small
\begin{tabular}{@{}lll@{}}
\toprule
Package & MaxEnt integration & Dependency \\
\midrule
\textbf{dismo} \citep{Hijmans2023}
  & rJava bridge to \texttt{maxent.jar} & Java JDK, rJava \\
\textbf{kuenm} \citep{Cobos2019}
  & \texttt{system2("java -jar maxent.jar")} & Java JDK \\
\textbf{kuenm2} (Cobos et al.)
  & \texttt{glmnet} via forked \texttt{maxnet} code & glmnet \\
\textbf{wallace} \citep{Kass2018wallace}
  & Delegates to ENMeval $\to$ maxnet or dismo & ENMeval \\
\textbf{ENMTools} \citep{Warren2010}
  & \texttt{dismo::maxent()} directly & dismo, rJava \\
\textbf{biomod2} \citep{Thuiller2009}
  & Java (\texttt{system2}) or \texttt{maxnet} & Optional Java or maxnet \\
\textbf{maxentcpp}
  & Native C++17 via Rcpp & None (self-contained) \\
\bottomrule
\end{tabular}
\end{table}

These packages can be grouped by their relationship to the MaxEnt
algorithm:

\begin{itemize}[nosep,leftmargin=*]
  \item \textbf{Black-box wrappers} (dismo, kuenm, ENMTools, biomod2
    MAXENT mode): invoke the Java binary and read output files.
    The optimizer is inaccessible.
  \item \textbf{Approximate reimplementations} (maxnet, kuenm2, biomod2
    MAXNET mode): use \texttt{glmnet} coordinate descent on the logistic
    regression formulation.
    The statistical model is equivalent but the optimization path
    differs, which can produce numerically different results for
    threshold and hinge features.
  \item \textbf{Faithful reimplementation} (maxentcpp): ports the
    original Sequential optimizer to C++ and proves per-iteration
    numerical equivalence.
\end{itemize}

The unique contribution of \texttt{maxentcpp} is that it is the only
package offering a compiled native reimplementation of the actual MaxEnt
optimizer---preserving both the statistical model and the optimization
algorithm while eliminating the Java dependency.

% ══════════════════════════════════════════════════════════════════════════
\section{Limitations and inherited assumptions}
% ══════════════════════════════════════════════════════════════════════════

The Maxent~3.4.4 algorithm that \texttt{maxentcpp} faithfully reproduces
has well-documented limitations that users should be aware of:

\begin{itemize}[nosep,leftmargin=*]
  \item \textbf{Feature explosion.}
    The number of features (particularly hinge and threshold) grows with
    the number of environmental variables, which can cause
    identifiability issues in high-dimensional settings
    \citep{Elith2011,Merow2013}.

  \item \textbf{Sensitivity to background sampling.}
    Maxent's output is influenced by the choice and size of the
    background sample, which affects the estimated density and can bias
    predictions in undersampled regions \citep{Phillips2009}.

  \item \textbf{$\ell_1$ regularization limitations.}
    The sequential coordinate ascent with $\ell_1$ penalties produces
    sparse models but does not guarantee selection of the ``correct''
    features under high correlation among predictors
    \citep{Hastie2015}.

  \item \textbf{No principled uncertainty quantification.}
    Unlike Bayesian approaches or bootstrap-based methods, the Maxent
    optimizer produces point estimates without confidence intervals.
\end{itemize}

A principled alternative would be to reformulate the Maxent objective
using modern convex optimization tooling (e.g., proximal gradient
methods, ADMM) or Bayesian inference.
However, such reformulations would produce different results than the
widely used Java implementation, breaking comparability with published
analyses.
The design choice in \texttt{maxentcpp} is explicitly to preserve this
comparability.

% ══════════════════════════════════════════════════════════════════════════
\section{Research impact statement}
% ══════════════════════════════════════════════════════════════════════════

\texttt{maxentcpp} is designed as a numerically faithful alternative to
Java Maxent for R-based ecological workflows.
The term ``drop-in replacement'' applies specifically to the implemented
features listed in the completeness table above: for those features,
\texttt{maxentcpp} produces results identical to Java Maxent (to within
IEEE~754 double-precision limits) given the same inputs, defaults, and
random seed.
Features not yet implemented (categorical variables, replicate runs,
jackknife) require continued use of Java Maxent.

The package is distributed under the MIT license (compatible with
Eigen's MPL2 and RcppEigen's GPL-2 via the ``or later'' clause), with
full source code, a pkgdown documentation website
(\url{https://alrobles.github.io/maxentcpp/}), and four vignettes
covering a getting-started walkthrough, the mathematical foundations of
Maxent features, a comparative motivation document, and a virtual
species validation study.
Continuous integration via GitHub Actions runs \texttt{R CMD check} on
Ubuntu, macOS, and Windows across R~release and R-devel.
The test suite comprises over 20~test files ($\sim$4\,800~lines)
covering feature generation, optimization convergence, spatial
projection, Java numerical equivalency, model diagnostics, and
end-to-end workflows.

By providing a dependency-free, memory-efficient, and numerically
faithful Maxent implementation, \texttt{maxentcpp} lowers the barrier
for large-scale biodiversity assessments, climate-change projections,
and conservation prioritization studies that rely on presence-only
species distribution models.

% ══════════════════════════════════════════════════════════════════════════
\section{AI usage disclosure}
% ══════════════════════════════════════════════════════════════════════════

Generative AI tools were used during the development of
\texttt{maxentcpp} and the preparation of this manuscript, as detailed
below.

\paragraph{Tools used.}
\begin{itemize}[leftmargin=*,nosep]
  \item \textbf{GitHub Copilot} (GPT-4-based, 2025 version): code
    scaffolding and boilerplate generation during the initial porting
    phases in \texttt{alrobles/Maxent} and
    \texttt{alrobles/maxentcpp-devel}.
  \item \textbf{Devin} (Cognition AI, 2025--2026): incremental feature
    implementation, test writing, documentation generation, CRAN
    compliance fixes, CI configuration, and manuscript drafting.
  \item \textbf{Claude} (Anthropic, 2025--2026): complex refactoring,
    cross-cutting documentation, and code review.
\end{itemize}

\paragraph{Nature and scope of assistance.}
AI tools contributed to: C++ code generation and refactoring from Java
source, Rcpp binding scaffolding, R wrapper functions, \texttt{testthat}
test scaffolding and expansion, \texttt{roxygen2} documentation,
vignette drafting, pkgdown site configuration, CI workflow setup, CRAN
compliance review, and manuscript drafting.
The core algorithmic C++ classes (\texttt{Sequential},
\texttt{FeaturedSpace}, feature types) were translated from Java with AI
assistance but under direct human supervision: each class was mapped
line-by-line from the corresponding Java source, reviewed for
correctness against the Java reference, and validated through the
\texttt{maxentcppCompTest} trajectory comparison framework.
Approximately 60\% of the C++ algorithmic lines were initially drafted
by AI tools and subsequently reviewed and corrected by the human author;
the remaining 40\% were written directly by the author, particularly
the performance-critical optimizer inner loops and numerical edge cases.

\paragraph{Maintainability.}
The human author (\'A.\,L.\ Robles~Fern\'andez) can independently
explain and modify every algorithmically significant class.
As evidence: the \texttt{Sequential::run()} optimizer,
\texttt{good\_alpha()}, \texttt{delta\_loss\_bound()},
\texttt{newton\_step\_feature()}, and \texttt{do\_parallel\_update()}
methods are documented with line-by-line references to the corresponding
Java source (e.g., ``mirrors Sequential.java:294..342''), and the author
has independently debugged numerical discrepancies during the fidelity
testing process.

\paragraph{Confirmation of review.}
All AI-generated code, tests, documentation, and manuscript text were
reviewed, edited, and validated by the human author
(\'A.\,L.\ Robles~Fern\'andez), who made all core design decisions
including the algorithm translation strategy, API design, numerical
fidelity requirements, and the phased development architecture.
The full development history is publicly available in the four
repositories listed above.

% ══════════════════════════════════════════════════════════════════════════
\section*{Acknowledgements}
% ══════════════════════════════════════════════════════════════════════════

The author thanks Steven~J. Phillips, Robert~P. Anderson, and
Robert~E. Schapire for creating the original Maxent software and making
the source code publicly available under an MIT license.
The \texttt{Rcpp}, \texttt{RcppEigen}, and \texttt{terra} package
maintainers are gratefully acknowledged for providing the infrastructure
that makes \texttt{maxentcpp} possible.

% ══════════════════════════════════════════════════════════════════════════
\bibliographystyle{plainnat}
\bibliography{paper}
% ══════════════════════════════════════════════════════════════════════════

\end{document}
